\section{习题课 \#5（10月11日）}

\subsection{线性方程组的存在性与唯一性}

考虑线性方程组
\begin{equation}
 Ax=b,
 \qquad
 A\in\F^{m\times n},\quad x\in\F^n,\quad b\in\F^m.
\end{equation}

\begin{theorem}[相容性判别]
下列条件等价：
\begin{enumerate}
  \item 方程组 $Ax=b$ 有解；
  \item $b\in\Image A$；
  \item $\Span(A,b)=\Span A$；
  \item $\rank(A,b)=\rank A$。
\end{enumerate}
\end{theorem}

\begin{proof}
矩阵 $A$ 的列向量记为 $\alpha_1,\ldots,\alpha_n$。方程
\begin{equation}
 Ax=b
\end{equation}
等价于
\begin{equation}
 x_1\alpha_1+\cdots+x_n\alpha_n=b,
\end{equation}
所以有解当且仅当 $b$ 属于 $A$ 的列空间。把 $b$ 添作一列不改变列空间，等价于增广矩阵与系数矩阵秩相同。
\end{proof}

\begin{theorem}[唯一性判别]
若 $Ax=b$ 已知有解，则下列条件等价：
\begin{enumerate}
  \item $Ax=b$ 的解唯一；
  \item 齐次方程 $Ax=0$ 只有零解；
  \item $\Ker A=\{0\}$；
  \item $A$ 的列向量线性无关；
  \item $\rank A=n$。
\end{enumerate}
\end{theorem}

\begin{proof}
若 $x_1,x_2$ 都是解，则
\begin{equation}
 A(x_1-x_2)=0.
\end{equation}
因此非齐次方程的解是否唯一，恰由齐次方程是否存在非零解决定。其余等价关系来自秩--零化度公式
\begin{equation}
 \dim\Ker A=n-\rank A.
\end{equation}
\end{proof}

\begin{corollary}
设 $A\in\F^{m\times n}$。
\begin{enumerate}
  \item 若 $\rank A=m$，则对每个 $b\in\F^m$，方程 $Ax=b$ 都有解；
  \item 若 $\rank A=n$，则对每个 $b\in\F^m$，方程 $Ax=b$ 至多有一个解；
  \item 若 $A$ 同时满行秩和满列秩，则 $m=n$ 且 $A$ 可逆，对每个 $b$ 都有唯一解。
\end{enumerate}
\end{corollary}

\begin{remark}
Gauss--Jordan 消元给出同一判别的计算版本。阶梯形增广矩阵中出现
\begin{equation}
 (0,\ldots,0\mid c),\qquad c\neq0,
\end{equation}
时无解；若无此矛盾行，则方程相容。相容时，主元列数等于未知数个数则唯一，否则含自由变量并有无穷多解。
\end{remark}

\subsection{解集的仿射结构}

\begin{theorem}
假设 $Ax=b$ 相容，并固定一个特解 $x_0$。则全部解构成仿射子空间
\begin{equation}
 \set{x\in\F^n:Ax=b}=x_0+\Ker A
 =\set{x_0+z:z\in\Ker A}.
\end{equation}
\end{theorem}

\begin{proof}
若 $Ax=b$，则
\begin{equation}
 A(x-x_0)=0,
\end{equation}
故 $x-x_0\in\Ker A$。反之，若 $z\in\Ker A$，则
\begin{equation}
 A(x_0+z)=Ax_0+Az=b.
\end{equation}
\end{proof}

\begin{remark}
因此，非齐次方程的全部信息可以分成两部分。一个特解决定仿射子空间的位置，齐次方程的零空间决定其方向。相容方程组的自由参数个数为
\begin{equation}
 \dim\Ker A=n-\rank A.
\end{equation}
\end{remark}

\subsection{从主元列构造零空间的一组基}

设
\begin{equation}
 A=(\beta_1,\ldots,\beta_n),
 \qquad
 \rank A=r,
\end{equation}
并设 $\beta_1,\ldots,\beta_r$ 是一组极大线性无关列。对每个 $j=r+1,\ldots,n$，存在唯一系数 $c_{1j},\ldots,c_{rj}$ 使
\begin{equation}
 \beta_j=\sum_{i=1}^r c_{ij}\beta_i.
\end{equation}
令
\begin{equation}
 \eta_j=(-c_{1j},\ldots,-c_{rj},0,\ldots,0,
 \underset{j\text{ 位}}{1},0,\ldots,0)^T.
\end{equation}
则 $A\eta_j=0$。

\begin{proposition}
向量
\begin{equation}
 \eta_{r+1},\ldots,\eta_n
\end{equation}
构成 $\Ker A$ 的一组基。
\end{proposition}

\begin{proof}
每个 $\eta_j$ 都属于 $\Ker A$。它们在第 $r+1,\ldots,n$ 个坐标上形成单位矩阵，故线性无关。数量为 $n-r=\dim\Ker A$，因而是一组基。
\end{proof}

\begin{remark}
这正是把阶梯形方程中的自由变量逐个取为 $1$、其余自由变量取为 $0$ 所得到的基础解系。
\end{remark}

\subsection{行空间与零空间的正交关系}

在 $\F=\R$ 或 $\C$ 上使用标准内积。复数情形中转置均应替换为共轭转置。

\begin{theorem}[基本正交分解]
对 $A\in\R^{m\times n}$，有
\begin{equation}
 \Image A^T=(\Ker A)^\perp,
 \qquad
 \R^n=\Image A^T\oplus\Ker A.
\end{equation}
相应地，
\begin{equation}
 \Image A=(\Ker A^T)^\perp,
 \qquad
 \R^m=\Image A\oplus\Ker A^T.
\end{equation}
\end{theorem}

\begin{proof}
若 $y=A^Tu$ 且 $z\in\Ker A$，则
\begin{equation}
 y^Tz=u^TAz=0,
\end{equation}
所以 $\Image A^T\subseteq(\Ker A)^\perp$。两边维数均为
\begin{equation}
 \rank A=n-\dim\Ker A,
\end{equation}
故二者相等。标准内积正定，因此子空间与其正交补直和为全空间。
\end{proof}

\begin{remark}
在一般域上，形式 $x^Ty$ 未必正定，因而不能无条件写出上述直和。例如在 $\F_2^2$ 中，非零向量 $(1,1)^T$ 满足
\begin{equation}
 (1,1)(1,1)^T=0,
\end{equation}
一个子空间可能与自身有非零交。因此原稿中特别强调，正交直和结论依赖于实数或复数上的正定内积。
\end{remark}

\subsection{典型消元题}

\begin{exercise}
判断方程组
\begin{equation}
\begin{cases}
 3x_1-5x_2+2x_3+4x_4=2,\\
 7x_1-4x_2-x_3+3x_4=5,\\
 5x_1+7x_2-4x_3-6x_4=3
\end{cases}
\end{equation}
是否相容。
\end{exercise}

\begin{solution}
对增广矩阵作初等行变换，会出现系数部分全为零而常数项非零的一行，因此
\begin{equation}
 \rank(A,b)>\rank A.
\end{equation}
故方程组无解。
\end{solution}

\begin{exercise}
求解
\begin{equation}
\begin{cases}
 2x_1+5x_2-8x_3=8,\\
 4x_1+3x_2-9x_3=9,\\
 2x_1+3x_2-5x_3=7,\\
 x_1+8x_2-7x_3=12.
\end{cases}
\end{equation}
\end{exercise}

\begin{solution}
消元后有三个主元，且无矛盾行，故解唯一。回代得到
\begin{equation}
 (x_1,x_2,x_3)^T=(3,2,1)^T.
\end{equation}
直接代回四个方程也可验证。
\end{solution}

\begin{remark}[含参数方程组的统一处理]
原稿还对一个含参数的四元方程组逐项讨论。计算时应先在不除以可能为零的参数表达式的前提下化为阶梯形，再按使主元消失的参数值分类。每一类中分别比较 $\rank A$ 与 $\rank(A,b)$，相容时再由 $n-\rank A$ 确定自由变量数。该流程比在一开始直接除以参数更不易漏解。
\end{remark}

\subsection{有限个真子空间不能覆盖无限域上的有限维空间}

\begin{theorem}
设 $V$ 是无限域 $\F$ 上的有限维向量空间。若
\begin{equation}
 V=W_1\cup\cdots\cup W_s
\end{equation}
且每个 $W_i$ 都是线性子空间，则至少有一个 $W_i=V$。换言之，有限个真子空间不能覆盖 $V$。
\end{theorem}

\begin{proof}
先看二维情形。若 $W_1,\ldots,W_s$ 都是真子空间，则每个至多是一条过原点的直线。取线性无关向量 $u,v$。由于 $\F$ 无限，向量
\begin{equation}
 u+tv,\qquad t\in\F,
\end{equation}
给出无穷多条不同方向，而有限条直线至多包含有限个这样的方向，故不能覆盖平面。

一般维数可归纳证明。选取一个不包含在给定真子空间中的超平面，或将问题限制到适当二维仿射平面，即可归约到上述论证。
\end{proof}

\begin{remark}
无限域假设不可删去。有限域上的有限维空间本身只有有限多个向量，也可以写成有限个一维子空间的并。
\end{remark}

\subsection{两个秩结论}

\begin{proposition}[斜对称矩阵的秩为偶数]
设 $A\in\F^{n\times n}$ 满足
\begin{equation}
 A^T=-A,
\end{equation}
且 $\operatorname{char}\F\neq2$。则 $\rank A$ 为偶数。
\end{proposition}

\begin{proof}
设 $r=\rank A$。可选出一个非零的 $r$ 阶子式。对斜对称矩阵，还可通过同时交换行和列把一个非奇异主子矩阵取作
\begin{equation}
 A_r\in\F^{r\times r},
 \qquad A_r^T=-A_r.
\end{equation}
于是
\begin{equation}
 \det A_r=\det A_r^T=
 \det(-A_r)=(-1)^r\det A_r.
\end{equation}
由于 $\det A_r\neq0$ 且特征不为 $2$，必有 $(-1)^r=1$，故 $r$ 为偶数。
\end{proof}

\begin{proposition}[添一列时秩至多增加一]
对 $A\in\F^{m\times n}$ 与 $b\in\F^m$，有
\begin{equation}
 0\leq\rank(A,b)-\rank A\leq1.
\end{equation}
并且
\begin{equation}
 \rank(A,b)=\rank A
 \quad\Longleftrightarrow\quad
 Ax=b\text{ 有解}.
\end{equation}
\end{proposition}

\begin{proof}
添入一个列向量至多使列空间维数增加 $1$。秩不增加恰好意味着 $b$ 已在原列空间中。
\end{proof}
